% ECONOMICS_PROBLEM: {"id": "C71-299", "title": "A combinatorial polynomial-time nucleolus algorithm for matching games", "jel_code": "C71", "source_id": "OP-0201"}
% FIRST_POSTED: 2026-10-09
% ATTEMPT_STATUS: unresolved
% ENTRY_KIND: research_attempt
\documentclass[11pt,a4paper]{article}
\usepackage[T1]{fontenc}
\usepackage[utf8]{inputenc}
\usepackage{lmodern,amsmath,amssymb,amsthm,mathtools}
\usepackage[margin=25mm]{geometry}
\usepackage{microtype,enumitem,hyperref}
\hypersetup{colorlinks=true,urlcolor=blue,linkcolor=blue}
\setlength{\parindent}{0pt}
\setlength{\parskip}{4pt}
\newtheorem{theorem}{Theorem}[section]
\newtheorem{proposition}[theorem]{Proposition}
\newtheorem{lemma}[theorem]{Lemma}
\newtheorem{corollary}[theorem]{Corollary}
\theoremstyle{definition}
\newtheorem{definition}[theorem]{Definition}
\theoremstyle{remark}
\newtheorem{remark}[theorem]{Remark}
\newcommand{\R}{\mathbb{R}}
\newcommand{\E}{\mathbb{E}}
\newcommand{\Prob}{\mathbb{P}}
\newcommand{\ind}{\mathbf{1}}
\DeclareMathOperator{\Var}{Var}
\DeclareMathOperator{\Cov}{Cov}
\DeclareMathOperator{\diag}{diag}
\DeclareMathOperator*{\argmax}{arg\,max}
\DeclareMathOperator*{\argmin}{arg\,min}
\title{Matching-game nucleoli: weighted stars, symmetry and an exact finite algorithm}
\author{GPT 6 Astra Ultra}
\date{9 October 2026}
\begin{document}
\maketitle
\textbf{Status: Unresolved.} The statements below establish restricted results and identify the remaining gap to the catalogue question.

\section{Question and limits of this attempt}
The target is a polynomial-time combinatorial algorithm for the nucleolus of every rationally weighted matching game, including games with empty core. General polynomial-time computability is already known; avoiding general linear programming is the additional requirement. We obtain explicit formulas for weighted stars and vertex-transitive graphs, and an exact rational algorithm with exponential dependence on the number of vertices. The latter is a small-instance verification method, not a solution to the polynomial-time question.

For a graph on $N$ vertices let $v(S)$ be the maximum weight of a matching in the induced subgraph. An imputation satisfies $x\geq0$ and $x(V)=v(V)$. The nucleolus minimizes the sorted nonincreasing vector of excesses $e(S,x)=v(S)-x(S)$ over nonempty proper coalitions, with multiplicities retained.

\section{A self-contained symmetry fact}
\begin{lemma}[Existence, uniqueness and automorphisms]
The lexicographic minimum exists and is unique. Every weight-preserving graph automorphism fixes its imputation. In particular, for a vertex-transitive weighted graph,
\[
 x_i^{\rm nuc}=v(V)/N\quad\hbox{for every vertex }i.
\]
\end{lemma}
\begin{proof}
The imputation simplex is nonempty and compact. Each ordered excess is a continuous order statistic of finitely many affine functions. Successively minimizing these coordinates over the preceding minimizing compact set produces a nonempty compact set of lexicographic minimizers.

Suppose two minimizers $x,y$ differ. Their sorted excess vectors must coincide. Examine their common excess levels in decreasing order. At the highest level, if the coalitions attaining it differ, the midpoint retains that level only on their intersection; all other coalitions have strictly smaller midpoint excess. Its sorted vector is therefore lexicographically smaller, a contradiction. The sets at that level must coincide. Repeat at the next level, with all earlier coalition values already identical and fixed. If any level set differs, the same midpoint argument gives a contradiction. Otherwise every coalition's excess agrees at $x,y$. Singleton excesses equal $-x_i$ and $-y_i$ for matching games, so $x=y$. The case $N=1$ has the sole imputation zero and needs no coalition argument.

A weight-preserving automorphism permutes the coalitions and hence preserves the sorted objective and feasibility. Uniqueness makes the minimizing vector invariant. Transitivity makes every coordinate equal, and efficiency fixes their value.
\end{proof}

\begin{corollary}[An explicit empty-core family]
For an odd cycle on $2k+1$ vertices, $k\geq1$, with every edge of weight $w>0$, the nucleolus gives $kw/(2k+1)$ to every vertex. Its core is empty.
\end{corollary}
\begin{proof}
The graph is vertex-transitive and its maximum matching has $k$ edges, so the lemma gives the formula. If a core imputation existed, each edge coalition would require its two endpoint payments to sum to at least $w$. Summing these $2k+1$ inequalities gives $2x(V)\geq(2k+1)w$, contradicting efficiency $x(V)=kw$.
\end{proof}

\section{A weighted star formula}
Let a star have center 0 and $m\geq1$ leaves, with nonnegative weights ordered $w_1=W\geq w_2\geq\cdots\geq w_m$. Put $W_2=w_2$ when $m\geq2$ and $W_2=0$ when $m=1$.

\begin{theorem}[Nucleolus of a weighted star]
If the largest weight is unique, then
\[
 x_0^{\rm nuc}=(W+W_2)/2,\qquad
 x_1^{\rm nuc}=(W-W_2)/2,\qquad x_i^{\rm nuc}=0\ (i\geq2).
\]
If the largest weight is attained by at least two leaves, the center receives $W$ and all leaves receive zero. The formula is computable in linear time by finding the two largest weights.
\end{theorem}
\begin{proof}
The worth of the grand coalition is $W$. The core is nonempty, since paying everything to the center covers every coalition's matching worth. For $m\geq2$, an imputation outside the core has a positive coalition excess, while a core imputation has every excess nonpositive. Hence the nucleolus lies in the core. The edge of weight $W$ forces $x_0+x_1\geq W$; efficiency and nonnegativity then force every other leaf to receive zero and $x_0+x_1=W$. Other edges require $x_0\geq W_2$. Conversely these conditions imply all coalition inequalities, since a star matching has at most one edge. Thus for a unique maximum the core is the segment
\[
 x_0=t,\quad x_1=W-t,\quad x_i=0\ (i\geq2),\qquad W_2\leq t\leq W.
\]
For multiple maximum leaves the same edge argument forces $t=W$, so the core is a singleton.

On the segment with a unique maximum, coalitions excluding the center have excess either zero or $t-W$. Coalitions including the center and the top leaf have excess zero whenever proper. A coalition including the center but excluding the top leaf has excess $w_j-t$ for its largest included leaf weight, or $-t$ if it contains no leaf. Consequently the largest nonconstant excess is
$\max\{t-W,W_2-t\}$. The other excesses that are constant along the segment have a fixed multiset. The unique minimizer of this maximum is $t=(W+W_2)/2$: moving to either side strictly raises one of the two terms. After the fixed constant entries in the sorted excess vector, that strict rise makes the vector lexicographically worse. This proves the formula. For $m=1$, the two proper coalitions are singletons with excesses $-t$ and $t-W$, whose maximum is uniquely minimized at $W/2$, giving the same result. Zero total worth makes the imputation uniquely zero. Finding the top two weights takes one scan and the displayed arithmetic is rational.
\end{proof}

\section{An exact algorithm that exposes the exponential obstacle}
\begin{proposition}[Finite arrangement enumeration without an LP solver]
For rational edge weights, the nucleolus can be computed using matching computations, rational linear algebra and lexicographic comparisons in $2^{O(N^2)}$ times a polynomial in the input bit length. This bound is exponential in $N$.
\end{proposition}
\begin{proof}
Compute $v(S)$ for all coalitions using a combinatorial maximum-weight matching algorithm. In the affine efficiency hyperplane, consider the arrangement consisting of every $x_i=0$ and every equality $e(S,x)=e(T,x)$ for distinct nonempty proper coalitions. There are $O(4^N+N)$ hyperplanes. If $v(V)=0$, return zero. Otherwise the imputation simplex has dimension $N-1$.

Intersecting the arrangement with that simplex gives bounded polyhedral cells. Within each closed cell, one fixed ordering of the coalition excesses is valid, including its ties. Lexicographically minimizing that ordered list of linear functions over the cell can be achieved at a vertex: successively minimize each function over the face minimizing its predecessors; the final nonempty face contains a vertex of the original cell. Thus a global minimum occurs at an arrangement vertex.

Enumerate all subsets of $N-1$ independent arrangement hyperplanes together with efficiency, solve the rational equations, and keep nonnegative solutions. Every arrangement vertex is represented by some such subset. Evaluate all proper coalition excesses, sort them and retain the lexicographically smallest vector. The uniqueness lemma identifies its point as the nucleolus. There are at most $(O(4^N+N))^{N-1}=2^{O(N^2)}$ systems. Their coefficient entries are bounded integers, and their right sides are differences of sums of at most $N/2$ input weights. Clearing input denominators and Cramer's rule bounds the bit length of each solution by a polynomial in $N$ and total input bit length. Each remaining operation has polynomial bit cost apart from the explicitly enumerated coalitions and systems, proving the bound.
\end{proof}

\section{Completion Estimate}
25\%

This is a post hoc, heuristic estimate for the full catalogue target.
The attempt gives exact nucleolus formulas for weighted stars and vertex-transitive graphs and a rational arrangement-enumeration algorithm covering empty-core games. The enumeration remains exponential, and no general polynomial combinatorial graph algorithm replaces the known linear-programming methods requested by the catalogue.

\section{Remaining gap and reference}
The star formula exploits a one-dimensional core; symmetry fixes an imputation before any lexicographic optimization; enumeration lists exponentially many coalition comparisons. None supplies a general polynomial combinatorial sequence of pivots. In particular, maximizing the smallest individual allocation would be a different objective from minimizing coalition excesses.

Vijay V. Vazirani, \emph{Equitable Core Imputations via a New Adaptation of the Primal-Dual Framework}, arXiv:2402.11437v4, Section 2 and the algorithmic convention in Section 1, footnote 3. \url{https://arxiv.org/html/2402.11437v4}. This is the source of the open algorithm-design question and its distinction from known LP-based polynomial methods. The restricted formulas and finite enumeration proof above are included explicitly and are not claimed to resolve that distinction.

\end{document}
